% 来源：OEBPS/Text/01_23.xhtml（OPF spine 顺序；仅第一册）
\cleardoublepage
\BookEntry[加勒特·哈丁《生活在极限之内》]{加勒特·哈丁《生活在极限之内》}

\BookCover{images/书单10生活在极限之内.jpg}{[美] 加勒特·哈丁著| 戴星翼等译| 上海译文出版社2016年版}

\begin{BookSummaryBox}
\textbf{杰出生态哲学家哈丁指出，因为人类的自私，如果不遏制人口增长，地球的资源将会被过度消耗}
\par
\end{BookSummaryBox}

\subsubsection*{这本书谈什么？}

熟悉经济学的朋友都听说过著名的“公地悲剧”理论：以公共利益为代价的个人获利将会使得村庄的公共牧场被过度放牧，因为每个牧人都想放养更多的牲口，过度放牧的结果由所有牧人承担，而增加牲口的利益则为个体牧人所有。提出“公地悲剧”理论的人，正是这本书的作者，杰出的生态哲学家加勒特·哈丁。
\par

哈丁将“公地悲剧”理论运用于全球生态问题，认为长期以来的先污染一个地方然后再迁往他处的做法已经行不通了。人类的存在已经遍及地球，无处可逃。在哈丁看来，全世界仿佛置身于一只救生艇之中，而这只救生艇所能够容纳的人数是有上限的，超过就会沉没。传统上人们对无极限增长的乐观相信忽略了重要的一点：地球的承载能力是有限的。因此，控制人口数量刻不容缓。
\par

\subsubsection*{芒格为什么谈起这本书？}

1993年6月30日，在接受《杰出投资者文摘》的采访时，芒格说：
\par

\begin{quotation}
我推荐每个人都来读加勒特·哈丁的《生活在极限之内》这本书。这本书总结了哈丁的毕生所学。他是一个真正的思想家。这本书棒极了。
\par
\end{quotation}

在西科金融公司1997财年的年会上，芒格说：
\par

\begin{quotation}
写这本书（《生活在极限之内》）的时候，（加勒特·哈丁）已经78岁了。这本书我推荐你们读两遍（我就读了两遍），因为那是一个特别聪明的人的智慧结晶。
\par
\end{quotation}

\subsubsection*{《生活在极限之内》作者加勒特·哈丁如是说}

“人类认为将要发生的事能够影响所发生的事。”
\par

“技术进步已成为一种宗教……把财富带到商人门口的技术都被认为是进步。”
\par

“置于虚假之上的行动将导致灾难。”
\par

“环绕着复利机构的沉默的阴谋是易于理解的。为了鼓励工人的忠诚，那些掌管着社会经济系统的人感到，他们必须宣称这个系统是绝对稳定的。并且，如我们所知，一些银行家甚至有勇气将‘永久’一词体现在他们机构的名称中。(人们很容易想象出，对于一个背负着‘麻烦到来前的永久国家银行’这个最诚实的名字的机构，会发生什么。)残酷而不可否认的事实是，在一个物质资源有限的世界上——它是我们可获得的唯一世界——复利本身造成了内在不稳定的系统。”
\par
